\documentclass[11pt,a4paper]{article}
\usepackage[margin=24mm]{geometry}
\usepackage{xcolor}
\usepackage[mode=preprocess]{xetex-stack-renderer}
\definecolor{ink}{HTML}{18354B}
\definecolor{accent}{HTML}{267B82}
\setlength{\parindent}{0pt}
\setlength{\parskip}{10pt}
\newcommand\PageTitle[1]{{\color{ink}\LARGE\bfseries XSR 0.10 / Inline glyphs}\par
 {\color{accent}\large #1}\par\medskip\hrule\medskip}
\GlyphRegister{square}{square.png}
\GlyphRegister{tall}{tall.png}
\GlyphRegister{very-tall}{very-tall.png}
\GlyphRegister{wide}{wide.png}
\GlyphRegister{very-wide}{very-wide.png}
\GlyphRegister{vector}{square.svg}
\GlyphRegister{jpeg}{square.jpg}
\GlyphRegister{pdf}{square.pdf}
\begin{document}
\PageTitle{1. Named signs in running text}
An external sign is a text-sized visual glyph. A label points to an ordinary
file, and each occurrence follows the size of the surrounding type.

In this invented reading sample, a square sign \Glyph{square} separates two
phrases. The same sign \Glyph{square} occurs later in the sentence, and again
here \Glyph{square}. These are neutral stroke fixtures, not transcriptions of
historical characters.

\textbf{The complete public registration API}
\begin{verbatim}
\GlyphRegister{square}{square.png}
Before \Glyph{square} after.
\end{verbatim}

\textbf{A common square canvas, four source formats}

{\fontsize{26}{36}\selectfont
PNG \Glyph{square}\quad JPEG \Glyph{jpeg}\par
PDF \Glyph{pdf}\quad SVG \Glyph{vector}\par}

The three image assets use the same policy, including PDF page inclusion.
Their canvas is 1.10 times the local cell. SVG uses 1.00 and remains PGF paths;
PNG and JPEG remain embedded raster images. All four preserve their canvas
proportions and pass through the same inline layout model.

\textbf{Canvas margins are part of the supplied sign}

The original fixture leaves a small white border around its strokes. XSR keeps
that border. Prepare a reasonably cropped source before registering it: XSR
performs no background removal, thresholding, cropping, trimming or tracing.
\vfill
\small This four-page document is preprocessed once and compiled with
\texttt{xelatex -no-shell-escape}. No font or external glyph dataset is bundled.
\newpage
\PageTitle{2. Rectangular signs retain their proportions}
The geometric mean of the final width and height is tied to the current
ideographic cell. A tall sign stays tall; a wide sign occupies more horizontal
space. No dimension is clamped to 1em.

\textbf{Aspect ratios are canvas width : canvas height}

{\fontsize{24}{32}\selectfont
1:1 \Glyph{square}\qquad 1:2 \Glyph{tall}\qquad 1:4 \Glyph{very-tall}\par
\medskip
2:1 \Glyph{wide}\qquad 4:1 \Glyph{very-wide}\par}

\textbf{Natural line spacing in a paragraph}

\begingroup\setlength{\parskip}{0pt}
\noindent An ordinary first line establishes the normal text baseline.\\
A second ordinary line follows at the normal baseline distance.\\
This line contains a tall sign \Glyph{very-tall} within running text.\\
The next line moves down to respect the sign's reported depth.\\
Normal line spacing resumes after the tall sign has passed.\par
\endgroup

\textbf{Horizontal occupancy}

Before \Glyph{square} after. Before \Glyph{wide} after.
Before \Glyph{very-wide} after.

The box reports the complete canvas height and depth. Image side bearings
reduce advance by two percent of canvas width on each side, allowing a small,
intentional overhang into the usual white margin. No pixels are clipped.
\vfill
\small For canvas ratio $r=W/H$, cell size $C$ and uniform scale $s$:
$w=\alpha Cs\sqrt{r}$ and $h=\alpha Cs/\sqrt{r}$.
\newpage
\PageTitle{3. Size, center alignment and two adjustments}
A cached glyph follows the active font size at each use. Image glyphs align
their geometric center to the current ideographic cell center, followed by a
small downward optical correction. User raise is an additional displacement.

{\small Small text: before \Glyph{square} and \Glyph{tall} after.\par}
Normal text: before \Glyph{square} and \Glyph{tall} after.\par
{\Large Large text: before \Glyph{square} and \Glyph{tall} after.\par}

\textbf{Uniform scale}

Original \Glyph{square}\quad
Smaller \Glyph[scale=.8]{square}\quad
Larger \Glyph[scale=1.5]{square}

\textbf{An additional vertical raise}

Default \Glyph{square}\quad
Raised \Glyph[raise=2pt]{square}\quad
Lowered \Glyph[raise=-2pt]{square}

\begin{verbatim}
\Glyph[scale=1.5,raise=2pt]{square}
\end{verbatim}

Scale multiplies both dimensions and the policy correction by the same factor.
Raise changes height and depth honestly. Width, height, xscale, yscale, stretch and crop are not public
adjustments. There is no preserve-aspect switch: proportions are always kept.

\textbf{Vector and image baseline behavior}

Vector \Glyph{vector}\quad Image \Glyph{square}\quad
Vector \Glyph{vector}\quad Image \Glyph{square}

Vectors use the cell baseline and normal side spacing. Images use center
alignment plus a downward bias of $0.12$ times the scaled cell, slightly larger
normalization and tighter side bearings.
When no ideographic metrics exist in the current font, the reference is its
em square, extending from the baseline to one em above it.
\vfill
\small The same source response is reused across all sizes and adjustments on
this page. Geometry is resolved before rendering; current type metrics remain local.
\newpage
\PageTitle{4. One layout pipeline, separate input paths}
Manual assets, direct SVG and externally produced SVG all resolve to the same
visual-asset model before sizing, alignment, spacing and TeX box construction.
Provider configuration is separate from the two-argument label registry.

\textbf{Registered SVG and direct SVG}

{\fontsize{26}{36}\selectfont
\Glyph{vector}\quad\xsrVectorGlyph{square.svg}\quad\Glyph{vector}\par}

The first and third signs use a registered label. The middle sign uses the
existing direct SVG API. Their layout policy and vector payload are identical.
The existing provider API feeds this same stage after obtaining and validating
an SVG file.

\textbf{Repeated names across formats}

A PNG sign \Glyph{square} can recur as \Glyph{square}, a JPEG sign
\Glyph{jpeg} as \Glyph{jpeg}, and a PDF sign \Glyph{pdf} as \Glyph{pdf}.
A repeated label must still refer to the same literal file. Unknown labels and
conflicting registrations produce explicit errors.

\textbf{The boundary of this release}

XSR reads canvas metadata and places the original image or validated vector.
This release does not implement image cleanup, an IDS engine, KAGE, a network
glyph service, transcription metadata or PDF semantic text handling.

The stroke fixtures are original drawings created by the showcase builder.
The source files remain in the temporary build directory; only this showcase
PDF and its reproducible source are committed.
\vfill
\small XSR 0.10 / Generic inline external glyphs / Four-page verification specimen
\end{document}
